\documentclass[a1paper,portrait,20pt]{tikzposter}

\usepackage[utf8]{inputenc}
\usepackage[T1]{fontenc}
% Use Latin Modern fonts to avoid missing font shapes warnings
\usepackage{lmodern}
% Improve font rendering and protrusion
\usepackage[final]{microtype}
\usepackage[english]{babel}
\usepackage{amsmath,amssymb}
\usepackage{graphicx}
\usepackage{booktabs}
\usepackage{enumitem}
\usepackage{wrapfig}

\definecolor{MainBlue}{HTML}{0077B6}
\definecolor{DarkBlue}{HTML}{023E8A}

\definebackgroundstyle{CleanWhite}{
    \fill[white] (bottomleft) rectangle (topright);
}

\defineblockstyle{MinimalBlock}{
    titlewidthscale=1.0,
    bodywidthscale=1.0,
    titleleft,
    titleoffsetx=0pt,
    titleoffsety=0pt,
    bodyoffsetx=0pt,
    bodyoffsety=0pt,
    bodyverticalshift=0pt,
    roundedcorners=0,
    linewidth=0pt,
    titleinnersep=2pt,
    bodyinnersep=5pt
}{
    \fill[white] (blockbody.south west) rectangle (blockbody.north east);
    \ifBlockHasTitle
        \fill[white] (blocktitle.south west) rectangle (blocktitle.north east);
        \draw[MainBlue, line width=2pt] (blocktitle.south west) -- (blocktitle.south east);
    \fi
}

\definetitlestyle{MinimalTitle}{
    width=0.9\paperwidth,
    roundedcorners=0,
    linewidth=0pt,
    innersep=5pt,
    titletotopverticalspace=5pt,
    titletoblockverticalspace=15pt
}{
    \fill[white] (\titleposleft,\titleposbottom) rectangle (\titleposright,\titlepostop);
    \draw[MainBlue, line width=3pt] (\titleposleft,\titleposbottom) -- (\titleposright,\titleposbottom);
}

% Custom title with logos stacked on left, title on right
\makeatletter
\renewcommand\TP@maketitle{%
    \noindent
    \parbox{8cm}{%
        \centering
        \IfFileExists{figures/logo.png}{%
            \includegraphics[height=5.25cm,width=7.5cm,keepaspectratio]{figures/logo.png}%
        }{%
            \fcolorbox{MainBlue}{white}{\parbox{5cm}{\centering\vspace{1cm}\textcolor{gray}{LOGO 1}\vspace{1cm}}}%
        }
        \IfFileExists{figures/logo2.png}{%
            \includegraphics[height=5.25cm,width=7.5cm,keepaspectratio]{figures/logo2.png}%
        }{%
            \fcolorbox{MainBlue}{white}{\parbox{5cm}{\centering\vspace{1cm}\textcolor{gray}{LOGO 2}\vspace{1cm}}}%
        }
    }\hfill
    \parbox{0.8\linewidth}{%
        \raggedleft
        {\bfseries\fontsize{50}{50}\selectfont\color{titlefgcolor}\@title\par}
        \vspace{0.5em}
        {\fontsize{32}{38}\selectfont\@author\par}
        \vspace{0.3em}
        {\fontsize{26}{32}\selectfont\@institute}
    }%
}
\makeatother

\usetheme{Default}
\usebackgroundstyle{CleanWhite}
\useblockstyle{MinimalBlock}
\usetitlestyle{MinimalTitle}
\tikzposterlatexaffectionproofoff

% Make block titles smaller
\makeatletter
\def\TP@blocktitlefont{\Large\bfseries}
\makeatother

\colorlet{titlefgcolor}{DarkBlue}
\colorlet{blocktitlefgcolor}{DarkBlue}
\colorlet{blockbodyfgcolor}{black}

% Pale blue color for theory blocks
\definecolor{TheoryBlue}{HTML}{BBDEFB}

% Define a block style for theoretical content: single rounded rectangle for title+body
\defineblockstyle{TheoryBlock}{
    titlewidthscale=1.0,
    bodywidthscale=1.0,
    titleleft,
    titleoffsetx=0pt,
    titleoffsety=0pt,
    bodyoffsetx=0pt,
    bodyoffsety=0pt,
    bodyverticalshift=0pt,
    roundedcorners=10,
    linewidth=0pt,
    titleinnersep=5pt,
    bodyinnersep=10pt
}{
    % Single rounded rectangle covering both title and body
    \ifBlockHasTitle
        \filldraw[rounded corners=10pt, fill=TheoryBlue, fill opacity=0.45, draw=MainBlue!70, line width=2pt]
            (blockbody.south west) rectangle (blocktitle.north east);
        % Underline under title to separate it from body
        \draw[MainBlue!80, line width=1.5pt]
            (blocktitle.south west) -- (blocktitle.south east);
    \else
        \filldraw[rounded corners=10pt, fill=TheoryBlue, fill opacity=0.45, draw=MainBlue!70, line width=2pt]
            (blockbody.south west) rectangle (blockbody.north east);
    \fi
}

% Convenience macro to use the theory block style for specific blocks
% Must restore MinimalBlock after use to prevent style leakage
\newcommand{\theoryblock}[2]{%
    \useblockstyle{TheoryBlock}%
    \block{#1}{#2}%
    \useblockstyle{MinimalBlock}%
}

% Define a block style with negative vertical shift to overlap with previous block
\defineblockstyle{OverlapBlock}{
    titlewidthscale=1.0,
    bodywidthscale=1.0,
    titleleft,
    titleoffsetx=0pt,
    titleoffsety=0pt,
    bodyoffsetx=0pt,
    bodyoffsety=-6cm,
    bodyverticalshift=-6cm,
    roundedcorners=0,
    linewidth=0pt,
    titleinnersep=2pt,
    bodyinnersep=5pt
}{
    \fill[white] (blockbody.south west) rectangle (blockbody.north east);
    \ifBlockHasTitle
        \fill[white] (blocktitle.south west) rectangle (blocktitle.north east);
        \draw[MainBlue, line width=2pt] (blocktitle.south west) -- (blocktitle.south east);
    \fi
}

\title{\textbf{Prisoner's Dilemma Generalizations \& NLHF}}
\author{Nikolai Vinogradov, Arsenii Burzilov, Ignashin Igor, Krylov Alexey}
\institute{MIPT}

\begin{document}
\maketitle

\begin{columns}
\column{0.50}

\block{Game Model (Prisoner's Dilemma)}{
\textit{Classic game where 2 players choose to cooperate or defect.}

Each agent action: $a\in[0,1]$ (cooperation degree).

\textcolor{MainBlue}{\textbf{Payoff:}} $T$=temptation, $R$=reward, $P$=punishment, $S$=sucker

$R(a_i,a_j)=R a_i a_j + T(1-a_i)a_j + S a_i(1-a_j) + P(1-a_i)(1-a_j)$
}

\theoryblock{SARSA + Boltzmann Policy}{
\textit{SARSA: on-policy RL algorithm learning Q-values from $(s,a,r,s',a')$ tuples.}

\textcolor{MainBlue}{\textbf{Q-function:}} expected cumulative reward for action $a$ in state $s$

$Q(s_t,a_t) \leftarrow Q + \alpha(r_t + \gamma Q(s_{t+1},a_{t+1}) - Q)$

\textcolor{MainBlue}{\textbf{Boltzmann:}} soft action selection ($\beta$=inverse temperature)

$\pi(a)=\frac{e^{\beta Q(a)}}{\sum_{a'} e^{\beta Q(a')}}$
}

\theoryblock{Stationary Distribution}{
\textit{$\gamma$: discount factor (future reward weight), $\alpha$: learning rate}

At convergence: $\pi^*(a) = \frac{e^{\beta Q^*(a)}}{\sum_{a'} e^{\beta Q^*(a')}}$

Bellman equation: $Q^*(a) = \mathbb{E}[r(a,a') + \gamma Q^*(a')]$

\textbf{Key:} $\pi^*(a)$ independent of $\gamma$, only convergence speed varies. Also we could say that about Q-learning with argmax.
}

\block{SARSA: Two-Player Results}{
Graphic for $R$ = 3, $T$ = 5, $P$ = 0, $S$ = 1, $\beta$ = 2.0, $\gamma$ = 0.6, steps = 2M

\begin{center}
\includegraphics[width=0.7\linewidth]{beta2.0_gamma0.6.png}
\end{center}
}

\block{Q-Learning \& Results: Multi-Player (N=5)}{
\textit{High $\beta$: escaping strategic trap causes oscillations.}

\begin{center}
\includegraphics[width=0.60\linewidth]{multiplayer_oscillations.jpg}
\end{center}

\vspace{-0.5cm}

\textcolor{MainBlue}{\textbf{Continuous \& Multi-Player Extensions:}}
Continuous action space $[0,1]$ enables agents to learn nuanced mixed strategies rather than pure cooperate/defect. Multi-player games (N=5) exhibit complex phase-shifted oscillations between cooperation and defection states. Strategic trap: agents initially cooperate but eventually discover defection yields higher payoff, triggering collective shift. Oscillation mechanism: when all defect, payoffs drop; some agents rediscover cooperation, restarting the cycle.

\textcolor{MainBlue}{\textbf{Key Findings:}}
Higher $\beta$ (lower temperature) leads to sharper policy distributions and faster convergence toward defection equilibrium. Discount factor $\gamma$ affects convergence speed but not the final stationary equilibrium distribution.
}


\column{0.50}

\theoryblock{NLHF: Reward Model with KL Regularization}{
\textit{Train reward model from preferences with regularization to reference policy $\pi_0$.}

\textcolor{MainBlue}{\textbf{Bradley-Terry preference model:}} 
$P(y_w \succ y_l | x) = \sigma(r_\theta(x, y_w) - r_\theta(x, y_l))$

where $\sigma$ is the sigmoid function, $r_\theta$ is the learned reward model

\textcolor{MainBlue}{\textbf{NLHF Loss with KL regularization:}} \\
$\mathcal{L} = -\log\sigma\left(r_w - r_l + \tau \cdot \left(\log\pi_0(y_w|x) - \log\pi_0(y_l|x)\right)\right)$

where $r_w$, $r_l$ = rewards for chosen/rejected, $\tau$=KL weight

\textcolor{MainBlue}{\textbf{KL penalty prevents reward hacking:}} 
$D_{KL}(\pi \| \pi_0) = \mathbb{E}_{y \sim \pi}\left[\log\frac{\pi(y|x)}{\pi_0(y|x)}\right]$

Regularization term anchors learned rewards to base model preferences, preventing over-optimization
}

\block{NLHF: Implementation Pipeline}{
\textit{SFT: Supervised Fine-Tuning. LoRA: Low-Rank Adaptation.}

\textcolor{MainBlue}{\textbf{Steps:}}
\begin{enumerate}[leftmargin=1em,itemsep=0pt]
\item SFT on TL;DR dataset (15k samples, 3 epochs)
\item Create $N$=5 diverse policies via weight perturbation
\item Generate 100 summaries/policy (batched, validated)
\item Form preference pairs: heuristic scoring or human labels
\item Compute reference log-probs $\log\pi_0(y)$ for KL term
\item Train Reward Model (LoRA r=8) with NLHF loss
\end{enumerate}

\textcolor{MainBlue}{\textbf{Model:}} Qwen2.5-3B-Instruct, LoRA, $\tau$=0.1, lr=2e-5
}

\block{NLHF: Quality \& Results}{
\textcolor{MainBlue}{\textbf{Generation quality controls:}}
\begin{itemize}[leftmargin=1em,itemsep=0pt]
\item Multi-factor scoring: length, coherence, diversity
\item Garbage detection: repetition, non-printable chars
\item Validation with retry (2 attempts per prompt)
\item Batched inference with gradient checkpointing
\end{itemize}

\textcolor{MainBlue}{\textbf{Data options:}}
\begin{itemize}[leftmargin=1em,itemsep=0pt]
\item Synthetic pairs (heuristic): accuracy $\sim$60\%
\item OpenAI human feedback: accuracy $\sim$80\%
\end{itemize}

\textcolor{MainBlue}{\textbf{Memory optimizations:}}
Sequential forward passes, reduced batch size (2), gradient accumulation (32)
}

\block{NLHF: Training Results}{
\textcolor{MainBlue}{\textbf{Reward Distribution Analysis:}} Mean reward = 0.331, diverse responses per prompt

\begin{center}
\includegraphics[width=0.7\linewidth]{sample_rewards_distribution.png}
\end{center}
}

\theoryblock{Future Work: Nash-MD \& Game Theory}{
\textit{Current: Reward Model training. Next: Multi-agent equilibrium learning.}

\begin{wrapfigure}{r}{6.5cm}
\vspace{-1em}
\centering
\includegraphics[width=6cm]{qr-code-3.png}
\vspace{-0.5em}
{\small \textit{GitHub Repository}}
\vspace{-1em}
\end{wrapfigure}

\textcolor{MainBlue}{\textbf{Nash-MD (Mirror Descent):}} iteratively update policy toward Nash equilibrium

$\pi_{t+1} = \arg\max_\pi \left[\eta P(\pi \succ \pi_t^\mu) - \text{KL}(\pi \| \pi_t^\mu)\right]$

where $\pi_t^\mu = (1-\eta\tau)\pi_t + \eta\tau\mu$ (mixture with reference $\mu$)

\textcolor{MainBlue}{\textbf{Nash equilibrium:}} no player benefits from unilateral deviation

$\mathbb{E}_{y\sim\pi^*}[P(y \succ y')] \geq \frac{1}{2}, \forall \pi'$

\textcolor{MainBlue}{\textbf{Convergence guarantee:}} $O(\log T/\sqrt{T})$ to approximate Nash

\textcolor{MainBlue}{\textbf{Open questions:}} Multi-player dynamics, stability analysis
}

\end{columns}

\end{document}
